% ============================================================
% 6 讨论与结论
% ============================================================
\section{讨论与结论}
\label{sec:conclusion}

\subsection{讨论}

\chead{任务存在两个难度层级。}在干净、闭集的设定下，各方法在视频级都接近 $99\%$，
帧级则都在 $60\%$ 左右，方法之间难以区分；一旦引入几何扰动，差距立刻显现，
且裁剪会同时破坏检索与拒识（第 5.4、5.5 节）。光度退化则几乎不影响任何方法。
这说明该任务的判别信号主要来自低频背景的\textbf{空间布局}：
只削弱高频的扰动不构成考验，而改变布局的扰动才是。对评测而言，
报告干净图像上的 Top-5 会系统性高估方法在真实转存$/$裁剪场景下的能力。

\chead{帧级失效可以由建模解释，也可以由建模修复。}第 5.3 节的问题—假设—证据链条指出，
帧级止步于 $60\%$ 并非容量不足，而是建模在削弱微小空间信号。三条途径的检验结果并不相同：
全局池化与几何增强被证实确有伤害（去掉池化升 $12$ 个百分点、加入几何增强反而下降）；
在全局向量之外补一条逐位置监督只带来轮次波动量级的差别（第 5.3 节），
说明瓶颈不在“编码器看不到细节”，而在\textbf{检索时只用全局内积做比较}。
据此把整图比对换成更细粒度网格上的成对局部几何比较后，干净档帧级由 $60.00\%$ 升至 $67.67\%$
（第 5.6 节）。这一提升并非来自更强的描述子——初筛本身已是本协议下最强的一级方案——
也不是来自更多分辨率（零训练扫描中 $8\times8$ 即为峰值），
而是来自比较形式：候选之间逐位置校验局部是否对得上。

\chead{几何验证有两种同等可行的路线。}软互最近邻、可微 DLT 与软内点占比构成的重排器，
在几何七档上与 SIFT$+$RANSAC 相差约 $1$ 个百分点，且 $4\times4$ 参数网格表明这一差距
不能靠调参消除（第 5.8 节）。经典路线成熟稳定，可微路线则便于与检索端联合优化。
把几何一致性放进训练则不合适——无论作为辅助损失还是唯一目标都未奏效。

\chead{“帧级”必须按判据分档陈述。}第 5.7 节的稠密帧库实验改变了前几节“帧级”一词的含义。
当命中按帧相似度区间判定、且库中每条视频有 $100$ 张连续帧时，\code{clean} 档已有 $94.1\%$
（$\delta$ 取 $p90$）落在区间内，全局单路更高（$98.9\%$）；即使把区间收紧到 $p50$，
全局单路仍有 $85.2\%$，而几何重排为 $65.9\%$。
几何重排的优势只出现在几何退化档：\code{crop0.8} 上 $T1_{\text{loc}}$ 由
$5.6\%/19.1\%/32.6\%$ 升至 $19.4\%/51.3\%/68.1\%$，\code{crop0.7off} 由
$3.5\%/11.3\%/20.9\%$ 升至 $12.4\%/40.2\%/51.9\%$，视频级命中由 $62.2\%$ 升至 $74.1\%$
（三档依次为 $\delta=p50,p90,p95$）。
这与第 4.1 节 P4/P5 的预期一致：几何校验补的是退化下的局部对应，而不是干净帧之间的分辨力。
因此第 5.3、5.6 节的 $60\%\rightarrow67.67\%$ 只适用于“同视频相隔两秒的 $8$ 张里挑一张”这一口径，
不能外推为“在连续视频里定位”的增益。还要如实说明：门控与几何回退的阈值是在旧协议上标定的，
迁移到稠密协议后 \code{blur4} 的回退触发率达 $39.6\%$ 却伴随下降（$51.1\%$ 对全局单路的 $83.3\%$），
说明这些阈值需按判据重标定，本文未做。

\chead{局限。}实验主体基于单一数据集（FF++ C23 抽帧）与单一图像级任务。
在跨域泛化上，第 5.9 节用 Celeb-DF v2 做了\textbf{真实换脸}的跨数据集溯源（$589$ 库 / $680$ 查询，
命中背景源视频）：干净档 $98.82\%$，五档均值 $93.06\%$，对 SIFT$+$RANSAC 的 $91.50\%$ 领先 $1.56$ 个百分点，
几何裁剪与模糊档领先、\code{jpeg5} 略低 $0.29$ 个百分点；同协议的 real$\rightarrow$real 对照表明
换脸本身几乎不损失可溯源性，难度来自同视频不同帧与退化。
第 5.10 节另用 CelebAMask-HQ 构造了跨伪造类型（区域拼接）的补充压力测试，
在光度退化下复现了优势（\code{blur4} 领先 SIFT $17.3$ 个百分点），
但几何裁剪档与 SIFT 打平、绝对水平受“每源仅一张图库图”的协议限制，
且该集合为自建而非标准 benchmark。
两节都属补充证据：Celeb-DF 上的命中是\textbf{视频级}（背景源视频）而非精确帧，
更大规模的换脸语料（DFDC$/$DeeperForensics）与修复类伪造仍未展开。
HashVAE 未被成功复现，因此本文只引用其原文读数，未做同协议对照。
可微几何重排器当前依赖独立的局部描述子，其初筛召回限制了几何重排的上界。

\subsection{尚未解决的问题与下一步}

本文用检索期的成对比较补上了缺失的环节，但编码器本身仍未学会捕捉微小细节：
重排使用的是独立的局部描述子，初筛召回决定了它的上界。
第 5.3 节的两组单变量训练排除了两个解释：其一是负样本不够难（$58.99\%$ 对 $59.68\%$）；
其二是给全局向量补一条位置级监督（$8$ 选 $1$ 为 $60.14\%$ 对 $59.79\%$，
差值在轮次波动之内，且与零训练描述子的 $60.11\%$ 持平）。
后一组实验同时澄清了一个容易混淆之处：只要检索仍以单一全局向量的内积排序，
无论辅助损失怎么加，比较形式都没有改变。结合零训练分辨率扫描
（$8\times8$ 最好，$16\times16$ 以上递减，仅高频 $46.74\%$）可以排除“信息在更高分辨率里”的猜测。
因此剩余的核心解释是\textbf{检索时的比较形式}：全局内积只给出整体相似度，
无法校验局部是否对得上。第 5.6 节已按此取得 $66.25\%$ / $67.67\%$，
但其中仍有三个成分未被分离，下面三条路径都可以用现有 $8$ 选 $1$ 口径判定：

\begin{enumerate}[leftmargin=1.6em,itemsep=1pt]
\item \textbf{同一编码器的逐位置检索（无需训练）。}直接取第 5.3 节网格模型的 $64$ 个位置描述子，
      在候选之间做逐位置匹配并聚合打分，替换掉全局内积。若 $8$ 选 $1$ 显著高于 $60.14\%$，
      则证明瓶颈确实在“全局内积”而不在编码器的特征；若仍停在 $60\%$，
      则说明该特征本身已无更多可用信息。
\item \textbf{分离分辨率与显式对应。}把第 5.6 节 $32\times32$ 的局部描述子场降采样到 $8\times8$，
      用同一几何模块重排。若仍接近 $66\%$，说明收益来自对应关系；若掉回 $60\%$，
      说明分辨率本身也在起作用。
\item \textbf{冻结早期层，只训嵌入头。}若“训练后期梯度侵蚀空间布局”成立，
      那么在去掉池化、保留布局的基础上冻结前几个阶段，读数应能稳定在 $58\%$ 以上而不再回落；
      若仍然回落，则该解释不成立。
\end{enumerate}

另有一条可叠加的路径是\textbf{高频残差输入}：以原图减去高斯模糊版作为旁路分支，
放大同视频不同帧之间的细小差异。需要注意的是，零训练的高频描述子单独只有 $46.74\%$，
它更可能作为布局信息的补充而非替代。

此外，拒识与检索目前分别报告，二者的联合指标（先拒识再检索）尚未测量；
开集条件下的检索精度仍是本文框架已经支持但尚未展开的部分
（跨数据集泛化已由第 5.9、5.10 节部分补上）。

\subsection{结论}

本文围绕“以背景为来源指纹”的深度伪造溯源任务，构建了一套可复现的评测框架：
形式化了 \aura{} 任务与帧级$/$视频级判据，给出光度六档与几何六档的确定性退化基准，
引入闭集$/$开集拒识协议。在此之上，本文给出了一条完整的因果链——
帧级失效源于建模把微小空间信号交给单一全局向量的内积比对
（全局池化与几何增强进一步削弱它），而仅靠更难的负样本或给全局向量补逐位置监督
都无法越过零训练水平；把整图比对改为细粒度局部网格上的成对几何比较后，
干净档帧级由 $60.00\%$ 升至 $67.67\%$，几何七档均值与经典 SIFT$+$RANSAC 相当。
稠密帧库上的真帧级定位实验进一步限定了这一结论：干净与光度退化下全局外观已足以定位，
几何重排的增益集中在几何退化档（第 5.7 节）。
我们希望这套协议、这条因果链与其中的实测读数，为后续溯源研究提供一个可靠的对照起点。
